\documentclass[11pt]{article}
\usepackage[utf8]{inputenc}
\usepackage[T1]{fontenc}
\usepackage[a4paper,margin=1in]{geometry}
\usepackage{lmodern}
\usepackage{microtype}
\usepackage{amsmath,amssymb,mathtools}
\usepackage{mathrsfs}
\usepackage{tikz-cd}
\usetikzlibrary{arrows}
\usepackage{enumitem}
\usepackage{longtable,booktabs,array}
\usepackage[hidelinks]{hyperref}
\setlength{\parindent}{0pt}
\setlength{\parskip}{0.55em}
\emergencystretch=3em
\hbadness=10000
\hfuzz=20pt
\usepackage[normalem]{ulem}
\newcommand{\K}{K}
\newcommand{\Gr}{\operatorname{Gr}}
\newcommand{\Fil}{\operatorname{Fil}}
\newcommand{\ch}{\operatorname{ch}}
\newcommand{\Todd}{\operatorname{Todd}}
\newcommand{\cl}{\operatorname{cl}}
\newcommand{\Parf}{\operatorname{Parf}}
\newcommand{\Tor}{\operatorname{Tor}}
\newcommand{\Spec}{\operatorname{Spec}}
\newcommand{\RRR}{\mathrm{RRR}}
\providecommand{\codim}{\operatorname{codim}}
\providecommand{\Inv}{\operatorname{Inv}}
\providecommand{\Pic}{\operatorname{Pic}}
\providecommand{\Z}{\mathbb Z}
\providecommand{\Q}{\mathbb Q}
\providecommand{\Ahat}{\widehat A}
\providecommand{\Ahp}{\widehat A^{+}}
\providecommand{\That}{\widehat T}
\providecommand{\rad}{\operatorname{rad}}
\providecommand{\Gl}{\operatorname{Gl}}
\providecommand{\GL}{\operatorname{Gl}}
\providecommand{\ch}{\operatorname{ch}}
\providecommand{\Todd}{\operatorname{Todd}}
\providecommand{\Gr}{\operatorname{Gr}}
\providecommand{\red}{\text{réd}}

% Core macro additions
\providecommand{\Atil}{\widetilde A}
\providecommand{\Ctil}{\widetilde C}
\providecommand{\K}{\mathrm K}
\providecommand{\cC}{\mathcal C}
\providecommand{\cF}{\mathscr F}
\providecommand{\eps}{\varepsilon}
\providecommand{\ellc}{\ell}
\providecommand{\Tor}{\operatorname{Tor}}
\providecommand{\Coker}{\operatorname{Coker}}
\providecommand{\Ker}{\operatorname{Ker}}
\providecommand{\Supp}{\operatorname{Supp}}
\providecommand{\codim}{\operatorname{codim}}
\providecommand{\Gr}{\operatorname{G}}
\providecommand{\rank}{\operatorname{rank}}
\providecommand{\rang}{\operatorname{rang}}
\providecommand{\cO}{\mathcal O}
\providecommand{\sgacircled}[1]{\tikz[baseline=(char.base)]\node[shape=circle,draw,inner sep=.35pt](char){\mbox{\scriptsize #1}};}

\hypersetup{hypertexnames=false}
\newcommand{\sgastandalonefootnotetext}[2]{%
  \begingroup
  \NoHyper
  \footnotetext{#2}%
  \endNoHyper
  \endgroup
}
\newcommand{\sgarepeatfootnotemark}{%
  \begingroup
  \NoHyper
  \footnotemark[\value{footnote}]%
  \endNoHyper
  \endgroup
}
\newcommand{\sganohyperfootnotemark}{%
  \begingroup
  \NoHyper
  \footnotemark
  \endNoHyper
  \endgroup
}
\providecommand{\R}{\mathbb R}
\providecommand{\C}{\mathbb C}
\providecommand{\N}{\mathbb N}
\providecommand{\Zell}{\mathbb Z_\ell}
\providecommand{\Qell}{\mathbb Q_\ell}
\providecommand{\C}{\mathbb C}
\providecommand{\Ch}{\operatorname{Ch}}
\providecommand{\Sym}{\operatorname{Sym}}
\providecommand{\Td}{\operatorname{Td}}
% --- build fix (typesetting only), 15 Nov 2026 ---
\usepackage{adjustbox}
\setlength{\emergencystretch}{5em}
\tolerance=2000
\begin{document}
\begin{center}
{\Large Seminario de Geometría Algebraica del Bois Marie}\par
{\large 1966--67}\par
\vspace{1.5em}
{\Large Teoría de las intersecciones y teorema de Riemann--Roch}\par
{\large (SGA 6)}\par
\vspace{1.5em}
un Seminario dirigido por\par
P. Berthelot, A. Grothendieck, L. Illusie\par
con la colaboración de\par
D. Ferrand, J.-P. Jouanolou, O. Jussila, S. Kleiman, M. Raynaud, J.-P. Serre
\end{center}

\section*{Prefacio}
La presente edición del Seminario de Geometría Algebraica 6 reproduce casi
por completo la primera edición, distribuida por el I.H.E.S. en forma de
exposiciones mimeografiadas. Las correcciones efectuadas se refieren
esencialmente a erratas de imprenta, salvo en las Exposiciones I y III, donde
L. Illusie introdujo ligeras modificaciones en algunos enunciados incorrectos
de la versión primitiva. Por otra parte, se han añadido notas a pie de página,
en particular en la Exposición XIV por A. Grothendieck, con el fin de señalar
progresos recientes en cuestiones que todavía estaban abiertas durante el
Seminario.

Señalemos finalmente que la Exposición XI, de D. Ferrand, relativa a la teoría
del determinante para los complejos perfectos, no fue redactada y, por tanto,
no será publicada; para evitar errores de referencia, se ha conservado sin
cambios la numeración de las Exposiciones siguientes.

\begin{flushright}Junio de 1971\end{flushright}

\section*{Introducción}
El objeto del presente Seminario es desarrollar una teoría
\underline{global} de las intersecciones y una fórmula de Riemann--Roch para
esquemas arbitrarios. El lector encontrará una descripción del programa del
Seminario en la Exposición 0. La posibilidad, en principio, de demostrar una
fórmula de Riemann--Roch para esquemas \underline{regulares} generales,
siguiendo las líneas del conocido informe de Borel--Serre de 1958, estaba ya
clara en ese momento, al menos para un morfismo proyectivo. La extensión al
caso general es menos evidente; el programa en el que se inscribe (y que se
realiza parcialmente en nuestro Seminario) se remonta a 1960. Como ocurría
también con la teoría de dualidad de los haces coherentes (cf.\ R.~Hartshorne,
Residues and Duality, Lecture Notes in Mathematics
n\textsuperscript{o}~20, Springer), una herramienta esencial para formular
una teoría satisfactoria es la teoría de las \underline{categorías derivadas}
de Verdier, cuyo conocimiento es indispensable para comprender el Seminario.

Aparte de esta teoría, el estudio del Seminario apenas exige un conocimiento
general de los fundamentos de la Geometría Algebraica, tal como se exponen en
EGA, capítulos I, II y III; además, solo necesitaremos hacer referencias
ocasionales a EGA IV para ciertos hechos relativos a los morfismos planos o
lisos, que figuran también, en lo esencial, en las primeras Exposiciones de
SGA~1. Finalmente, para desarrollar ciertos resultados con toda la
generalidad deseable para las aplicaciones, utilizamos a veces la noción de
\underline{sitio anillado} y de \underline{topos anillado}, para la cual
remitimos a SGA~4, Exposiciones I a IV. El lector que desconozca el lenguaje
de los sitios y los topos podrá reemplazar en todas partes dichos objetos por
espacios topológicos ordinarios, reemplazando entonces los objetos del topos
por abiertos de estos espacios; no obstante, le aconsejamos preferentemente
que asimile el lenguaje de los topos, que proporciona un principio de
unificación sumamente cómodo.

La noción fundamental para la teoría aquí presentada es la de
\underline{complejo perfecto de módulos}, y sus diversas variantes,
desarrolladas en las Exposiciones I, II y III. Parece claro que estas nociones,
también importantes en otros contextos de la Geometría Algebraica (en
particular para las fórmulas de tipo Lefschetz--Verdier en cohomología con
coeficientes discretos (SGA 5) o coherentes), tendrán asimismo un papel que
desempeñar fuera de la Geometría Algebraica, en particular para formular una
variante analítica compleja del teorema de Riemann--Roch aquí presentado, o
variantes adecuadas del teorema del índice de Atiyah--Singer. En la Exp.\ II
se darán algunas indicaciones en este sentido. Desgraciadamente, todavía
falta en la actualidad un enunciado, siquiera heurístico, que englobe estos
dos teoremas (el primero de los cuales sigue siendo por ahora conjetural).
Parece faltar una idea nueva, como falta también en Geometría Algebraica para
llegar a una demostración del teorema de Riemann--Roch fuera de hipótesis
proyectivas (cf.\ Exp.\ XIV, n\textsuperscript{o}~2).

No haremos uso alguno en este Seminario de la \underline{teoría local} de las
intersecciones, expuesta en J.P. Serre, Algèbre Locale. Multiplicités (Lecture
Notes in Mathematics n\textsuperscript{o}~11, Springer), y nos limitaremos a
señalar aquí que este curso de Serre ejerció una influencia evidente en la
génesis de la teoría en 1957. Tampoco utilizaremos la teoría del
\underline{anillo de Chow} desarrollada en el Séminaire Chevalley 1958,
Exposiciones 2 y 3 (École Normale Supérieure). Esta teoría está vinculada de
manera esencial a hipótesis de no singularidad, mientras que el objeto de
nuestro Seminario es, por el contrario, desarrollar una teoría de las
intersecciones sobre esquemas generales (e incluso topoi localmente anillados
generales); véanse a este respecto los comentarios de la Exp.\ XIV,
n\textsuperscript{o}~4 y n\textsuperscript{o}~8, que exponen las relaciones
entre nuestra teoría y la del anillo de Chow y plantean la cuestión de una
generalización de esta última. Podemos decir que el Seminario presenta una
teoría de intersecciones coherente y autosuficiente, pero en modo alguno
exhaustiva respecto de los diversos puntos de vista útiles (e incluso
indispensables) en teoría de intersecciones, y que, por consiguiente, conviene
completarlo con las exposiciones citadas de Serre y Chevalley.

Se ha adjuntado al Seminario (como Apéndice a la Exp.\ 0) el informe de
Grothendieck de 1957 «~classes de Faisceaux et théorème de Riemann--Roch~»,
que había servido de base al seminario de Borel y Serre en Princeton ese mismo
año, así como a su artículo ya citado. Este informe esboza dos demostraciones
del teorema de Riemann--Roch para variedades cuasiproyectivas no singulares;
la primera de ellas, válida por el momento solo en característica cero, pero
que proporciona en cambio un resultado algo más preciso en el caso de una
inmersión, todavía no figura en la literatura (sin exceptuar el trabajo de
Borel--Serre ni el presente Seminario). Naturalmente, la lectura de este
informe no presupone la de ninguna otra Exposición del Seminario, e incluso
puede servir de introducción a su estudio al mismo título que la Exposición 0,
especialmente para un lector que aún no esté familiarizado con Borel--Serre.
La demostración a la que acabamos de aludir se aplica, además, esencialmente
sin cambios al caso de un morfismo proyectivo de espacios analíticos complejos,
y se aplicará sin duda también en característica arbitraria una vez resuelto el
problema de las operaciones «~potencias exteriores~» en la categoría derivada
(cf.\ Exp.\ XIV, n\textsuperscript{os}~1, 2, 3). La ausencia de un estudio de
esta cuestión constituye sin duda la laguna más grave de este Seminario, dentro
del propio enfoque que se ha adoptado en él.

Se observará la ausencia, en la tabla de materias del presente Seminario, de
dos de los participantes mencionados en la portada. Uno de ellos, J.P.
Jouanolou, había participado activamente en la elaboración técnica de la
primera parte del Seminario, pero no había podido participar en las
exposiciones orales; se le debe, en particular, la afirmación de independencia
lineal contenida en el importante enunciado VI 1.1, que da la estructura del
anillo $K^\circ$ de las clases de módulos localmente libres sobre un fibrado
proyectivo (que anteriormente solo se había demostrado bajo la hipótesis de
que el esquema base admitiese un módulo invertible amplio). El otro, J.P.
Serre, había impartido dos exposiciones orales, lógicamente independientes del
resto del Seminario, y por ello prefirió publicarlas en forma de artículo
separado (cf.\ J.P. Serre, Groupes de Grothendieck des schémas en groupes
réductifs déployés, de próxima aparición en Publications Mathématiques,
n\textsuperscript{o}~34).

Como en las demás partes del Seminario de Geometría Algebraica del Bois Marie,
las siglas SGA 1 a SGA 6 remiten a las distintas partes de dicho Seminario; el
número romano que sigue a estas siglas indica el n\textsuperscript{o} de la
Exposición, y las cifras arábigas que lo siguen corresponden a la numeración
interna de la Exposición de que se trate. Para las referencias internas al
presente Seminario se omiten las siglas SGA 6 y se utilizan referencias del
tipo I 4.9 (que significa: Exposición I, enunciado 4.9). La sigla EGA remite a
los Éléments de Géométrie Algébrique de J. Dieudonné y A. Grothendieck.

\section*{Tabla de materias}

\noindent
\underline{EXPOSICIÓN 0}\\
Esbozo de un programa para una teoría de las intersecciones sobre esquemas generales\\
por A. Grothendieck \dotfill 1\\[8pt]
Clases de haces y teorema de Riemann--Roch\\
por A. Grothendieck \dotfill 20\\[8pt]
\underline{EXPOSICIÓN I}\\
Generalidades sobre las condiciones de finitud en las categorías derivadas\\
por L. Illusie \dotfill 78\\[8pt]
\underline{EXPOSICIÓN II}\\
Existencia de resoluciones globales\\
por L. Illusie \dotfill 160\\[8pt]
\underline{EXPOSICIÓN III}\\
Condiciones de finitud relativas\\
por L. Illusie \dotfill 222\\[8pt]
\underline{EXPOSICIÓN IV}\\
Grupos de Grothendieck de los topos anillados\\
por L. Illusie \dotfill 274\\[8pt]
\underline{EXPOSICIÓN V}\\
Generalidades sobre los $\lambda$-anillos\\
por P. Berthelot \dotfill 297\\[8pt]
\underline{EXPOSICIÓN VI}\\
El $K^\circ$ de un fibrado proyectivo: cálculos y consecuencias\\
por P. Berthelot \dotfill 365\\[8pt]
\underline{EXPOSICIÓN VII}\\
Inmersiones regulares y cálculo del $K^\circ$ de un esquema obtenido por explosión\\
por P. Berthelot \dotfill 416\\[8pt]
\underline{EXPOSICIÓN VIII}\\
El teorema de Riemann--Roch\\
por P. Berthelot \dotfill 466\\[8pt]
\underline{EXPOSICIÓN IX}\\
Algunos cálculos de grupos $K_\bullet$\\
por P. Berthelot \dotfill 498\\[8pt]
\underline{EXPOSICIÓN X}\\
Formalismo de las intersecciones sobre los esquemas algebraicos propios\\
por O. Jussila\\
Con un Apéndice de A. Grothendieck\\
Especialización en teoría de intersecciones \dotfill 519\\[8pt]
\underline{EXPOSICIÓN XI} -- No redactada\\[8pt]
\underline{EXPOSICIÓN XII}\\
Un teorema de representabilidad relativa sobre el funtor de Picard\\
por M. Raynaud (redactado por S. Kleiman) \dotfill 595\\[8pt]
\underline{EXPOSICIÓN XIII}\\
Los teoremas de finitud para el funtor de Picard\\
por S. Kleiman \dotfill 616\\[8pt]
\underline{EXPOSICIÓN XIV}\\
Problemas abiertos en teoría de intersecciones\\
por A. Grothendieck \dotfill 667\\[8pt]
Índice terminológico \dotfill 691\\[4pt]
Índice de notaciones \dotfill 696\par
